\section{Conclusion}
\label{sec:conclusion}

We validated a structured LLM paper-scoring system against the public decisions
and reviewer ratings of a major machine learning venue, under a pre-registered,
fully reproducible protocol. The score separates rejected from accepted
submissions (AUROC~$\aurocMini$), rises monotonically across decision tiers, and
correlates with continuous reviewer ratings ($\rho=\spearmanRatingMini$). The
signal is strongest at the low end: weak submissions are consistently flagged,
with the lowest score quintile rejected at \bottomRejectRate\%, a
\bottomLift$\times$ lift over base rate, and strong work nearly absent from that
band. We confirmed the effect on a frontier-model cohort and showed that a cheap
model proxies the frontier score well enough to establish the relationship at
scale. Crucially, the signal is a product of the engineered pipeline rather than
the model it runs on: against a generic one-paragraph prompt on the same model and
PDF, the structured rubric-and-audit pipeline discriminates better and grades far
more consistently.

Pre-registering the claim, validating against two complementary ground truths,
generating every number from the data, and cataloguing the failure modes together
keep the result inside the boundary the evidence supports: an automated first pass that reliably
surfaces weak work for human attention, not a replacement for the judgment that
follows. That boundary follows the reliable-at-the-bottom asymmetry the human
process itself exhibits.

\paragraph{Reproducibility.}
The de-identified scored table, the decision and rating labels, the analysis
code, and the pre-registration are released; \texttt{run\_all.py} regenerates
every figure, table, and number in this paper from the data with a single
command. The grading rubric, prompts, and model configuration remain proprietary;
the manuscript-to-score function is audited through its outputs rather than
re-implemented.
